\documentclass[11pt]{article}

\usepackage[T1]{fontenc}
\usepackage{amsmath,amssymb,amsthm}
\usepackage[a4paper,margin=1in]{geometry}
\usepackage{microtype}
\usepackage[colorlinks=true,linkcolor=blue,citecolor=blue,urlcolor=blue]{hyperref}
\hypersetup{
  pdftitle={Maxwell Sharpness for Complements of Disjoint Triangles and Four-Cycles},
  pdfauthor={Xiaojian Zeng},
  pdfsubject={Generic rigidity and Maxwell dimension}
}

\newtheorem{theorem}{Theorem}[section]
\newtheorem{proposition}[theorem]{Proposition}
\newtheorem{lemma}[theorem]{Lemma}
\newtheorem{remark}[theorem]{Remark}

\newcommand{\R}{\mathbb{R}}
\newcommand{\K}{\mathbb{K}}
\newcommand{\Sym}{\operatorname{Sym}}
\newcommand{\rig}{\operatorname{rig}}
\newcommand{\rank}{\operatorname{rank}}
\newcommand{\coker}{\operatorname{coker}}

\title{Maxwell Sharpness for Complements of Disjoint Triangles and Four-Cycles}
\author{Xiaojian Zeng\\
  PBC School of Finance, Tsinghua University, Beijing 100190, China\\
  \texttt{xiaojianzeng@foxmail.com}}
\date{September 2026}

\begin{document}
\maketitle

\begin{abstract}
Let $D=sC_3\sqcup tC_4$ have $N=3s+4t\geq4$ vertices and let
$G=\overline D$.  If $a$ is the least nonnegative integer such that
$N\leq\binom{a+1}{2}$, then we prove
\[
 \rig(G)=M(G)=N-a-1,
\]
where $\rig$ and $M$ denote generic rigidity dimension and Maxwell
dimension.  The main ingredient is a symmetric-product theorem: for generic
$z_1,\ldots,z_N\in\K^a$, the tensors $z_i\odot z_j$ indexed by the edges of
$D$ are independent.  Its proof compares the exact compression codimension
of $Z^{\mathsf T}CZ=0$ with every rank stratum of the space of symmetric
matrices supported on the triangles and four-cycles.  An all-nonzero kernel
scaling then turns $Z$ into a Gale matrix, and stress-space duality shows that
deleting $E(D)$ from $K_N$ preserves the complete framework's rigidity rank.
This settles the pure exceptional-component family left outside the closest
bounded-degree-complement theorem.
\end{abstract}

\section{Statement and context}

A realization of a finite simple graph $H$ in $\R^d$ is \emph{generic} when
all of its coordinates are algebraically independent over $\mathbb Q$.
For $|V(H)|\geq d+1$, the theorem of Asimow and Roth identifies generic
rigidity with the rigidity-matrix rank condition
\[
 \rank R(H,p)=d|V(H)|-\binom{d+1}{2}.
 \tag{1.1}
\]
We use the convention that every graph is generically $0$-rigid.  The
rigidity dimension $\rig(H)$ is the greatest $0\leq d\leq |V(H)|-1$ for
which $H$ is generically $d$-rigid.  Maxwell's necessary edge count
\cite{Maxwell1864}, together with the generic rank criterion
\cite{AsimowRoth1978}, gives
\[
 \rig(H)\leq M(H),\qquad
 M(H):=\max\left\{0\leq d\leq |V(H)|-1:
 d|V(H)|-\binom{d+1}{2}\leq |E(H)|\right\}.
 \tag{1.2}
\]

Haslegrave, Michaeli and Nixon proved that the complement of a graph of
maximum degree at most two is Maxwell-sharp provided that no component of the
deleted graph is a triangle or a four-cycle
\cite[Theorem~1.4]{HaslegraveMichaeliNixon2026}.  They explicitly identified
the exclusions of $C_3$ and $C_4$ as limitations of their argument.  Their
Corollary~1.5 implies, for $D=sC_3\sqcup tC_4$, only
\[
 M(\overline D)-\rig(\overline D)
 \leq \left\lceil\frac{s+t}{2}\right\rceil,
 \tag{1.3}
\]
because one vertex must be deleted from each component of $D$ to remove all
triangle and four-cycle components.  The result below treats precisely this
pure excluded family.

\begin{theorem}\label{thm:main}
For all integers $s,t\geq0$ with $N=3s+4t\geq4$, let
$D=sC_3\sqcup tC_4$ and $G=\overline D$.  If
\[
 a=\min\left\{q\geq0:N\leq\binom{q+1}{2}\right\},
\]
then
\[
 \rig(G)=M(G)=N-a-1.
\]
\end{theorem}

The argument is self-contained apart from the standard generic-rigidity
criterion (1.1).  We work over an algebraically closed field $\K$ of
characteristic zero during the dimension calculation.  Since every open set
and polynomial used below is defined over $\mathbb Q$, the resulting
nonvanishing has a rational, and hence real, specialization.

\section{The Maxwell target}

The graph $D$ is $2$-regular, so it has $N$ edges.  Consequently
\[
 |E(G)|=\binom N2-N=\frac{N(N-3)}2.
 \tag{2.1}
\]

\begin{lemma}\label{lem:maxwell}
With $a$ as in Theorem~\ref{thm:main}, the Maxwell dimension of $G$ is
$N-a-1$.
\end{lemma}

\begin{proof}
Put $d=N-q-1$.  Substitution of (2.1) into the Maxwell inequality gives
\[
 dN-\binom{d+1}{2}\leq\frac{N(N-3)}2
 \quad\Longleftrightarrow\quad
 N\leq\frac{q(q+1)}2.
 \tag{2.2}
\]
Increasing $d$ by one decreases $q$ by one.  Thus the largest admissible
$d$ is obtained from the least $q=a$ satisfying the right side of (2.2),
and equals $N-a-1$.
\end{proof}

\section{Generic independence of the deleted-edge tensors}

For $x,y\in\K^a$, write
\[
 x\odot y:=\frac{xy^{\mathsf T}+yx^{\mathsf T}}2\in\Sym^2\K^a.
\]
The following proposition is the central algebraic statement.

\begin{proposition}[symmetric-product independence]\label{prop:compression}
Let $D=sC_3\sqcup tC_4$ have $N$ vertices.  If
$N\leq\binom{a+1}{2}$, then, for generic
$z_1,\ldots,z_N\in\K^a$, the $N$ tensors
\[
 \{z_i\odot z_j:ij\in E(D)\}
 \tag{3.1}
\]
are linearly independent in $\Sym^2\K^a$.
\end{proposition}

\begin{proof}
Let $Z$ be the $N\times a$ matrix with rows $z_i$, and let $L_D$ be the
$N$-dimensional space of symmetric $N\times N$ matrices whose diagonal
entries vanish and whose off-diagonal support is contained in $E(D)$.  A
linear relation among (3.1) is equivalent, up to a fixed nonzero factor on
each off-diagonal coefficient, to
\[
 Z^{\mathsf T}CZ=0
 \tag{3.2}
\]
for some nonzero $C\in L_D$.  We prove that the incidence variety of such
pairs $(C,Z)$ does not dominate the $aN$-dimensional space of $Z$.

\smallskip
\noindent\emph{Compression codimension.}
Fix a symmetric matrix $C$ of rank $r$.  Congruence reduces it to
$\operatorname{diag}(I_r,0)$.  After applying the same invertible row change
to $Z$, condition (3.2) becomes
\[
 X^{\mathsf T}X=0,
 \tag{3.3}
\]
where $X$ consists of the first $r$ rows.  On the locus $\rank X=k$, the
column space of $X$ is a totally isotropic $k$-plane in $\K^r$.  The variety
of these planes has dimension
\[
 k(r-k)-\binom{k+1}{2}.
\]
Indeed, the ordinary Grassmannian has dimension $k(r-k)$ and the vanishing
of the restricted nondegenerate symmetric form imposes the
$\binom{k+1}{2}$ independent entries of a symmetric Gram matrix.  Once the
plane is chosen, its spanning $r\times a$ matrices form an open subset of an
$ak$-dimensional vector space.  Thus the rank-$k$ part of (3.3) has dimension
\[
 f(k)=k(r-k)-\binom{k+1}{2}+ak.
\]
Here $0\leq k\leq\min(a,\lfloor r/2\rfloor)$, and
$f(k)-f(k-1)=a+r-3k+1>0$ throughout this range.  Therefore the codimension,
in the full space of $Z$, of (3.2) for this fixed $C$ is
\[
 c_a(r)=
 \begin{cases}
 \dfrac{h(2a-h+1)}2,&r=2h<2a,\\[5pt]
 a+\dfrac{h(2a-h-1)}2,&r=2h+1<2a,\\[5pt]
 \binom{a+1}{2},&r\geq2a.
 \end{cases}
 \tag{3.4}
\]

\smallskip
\noindent\emph{Rank strata in $L_D$.}
A weighted triangle block has the form
\[
 \begin{pmatrix}
 0&x&z\\ x&0&y\\ z&y&0
 \end{pmatrix}.
\]
Its determinant is $2xyz$, and its nonzero strata have
$(\text{rank},\text{dimension})=(2,2)$ or $(3,3)$.  Rank one is impossible.
After ordering a four-cycle by its bipartition, its block has the form
\[
 \begin{pmatrix}0&B\\B^{\mathsf T}&0\end{pmatrix},
 \qquad B\in\K^{2\times2},
\]
so its nonzero strata have $(\text{rank},\text{dimension})=(2,3)$ or
$(4,4)$.  Refining reducible strata into their finitely many irreducible
pieces does not change these dimension bounds.

Suppose that $u,v,x,y$ components are active in, respectively, the four
types
\[
 (2,2)_{C_3},\quad(3,3)_{C_3},\quad
 (2,3)_{C_4},\quad(4,4)_{C_4}.
\]
The corresponding product stratum has rank $r$ and dimension $\ell$ given
by
\[
 r=2u+3v+2x+4y,
 \qquad \ell=2u+3v+3x+4y=r+x.
 \tag{3.5}
\]
We claim that its projective dimension satisfies
\[
 \ell-1<c_a(r).
 \tag{3.6}
\]

If $r\geq2a$, then $\ell$ is at most the number of edge parameters in the
active components.  Hence
$\ell\leq N\leq\binom{a+1}{2}=c_a(r)$.

Let $r=2h<2a$.  Equation (3.5) gives $x\leq h$ and therefore
$\ell\leq3h$.  From (3.4),
\[
 c_a(2h)-(3h-1)=\frac{h(2a-h-5)+2}{2}.
 \tag{3.7}
\]
For $1\leq h\leq a-1$, the numerator is a concave quadratic in $h$, so
its minimum occurs at an endpoint.  The endpoint values are $2a-4$ and
$(a-2)(a-3)$.  They are nonnegative, with equality only when
$(a,h)=(2,1)$ or $(3,2)$.  Equality in $\ell\leq3h$ would require,
respectively, one or two active rank-two four-cycle blocks.  These require
$N\geq4$ or $N\geq8$, contradicting $N\leq\binom{a+1}{2}=3$ or $6$.
Thus (3.6) is strict also in the exceptional arithmetic cases.

Finally, let $r=2h+1<2a$.  Odd rank forces $v$ to be odd, hence at least
one rank-three triangle block is active.  It follows that $x\leq h-1$ and
$\ell\leq3h$.  Now
\[
 c_a(2h+1)-(3h-1)
 =\frac{2a+h(2a-h-7)+2}{2}>0.
 \tag{3.8}
\]
The numerator is again concave in $h$; at $h=1$ and $h=a-1$ its values are
$4a-6$ and $(a-2)(a-3)+2$, both positive.  This proves (3.6).

For the projectivization of any fixed rank-stratum piece, the incidence
pairs $([C],Z)$ satisfying (3.2) consequently have dimension at most
\[
 (\ell-1)+(aN-c_a(r))<aN.
\]
There are only finitely many choices of active blocks, block ranks and
irreducible pieces.  Their images lie in a proper closed subset of the space
of $Z$.  Thus a generic $Z$ admits no nonzero $C\in L_D$ satisfying (3.2),
which proves the proposition.
\end{proof}

\section{Gale normalization and rigidity}

We next pass from the unconstrained rows in
Proposition~\ref{prop:compression} to the Gale matrix of an affine
configuration.  Gale transforms and equilibrium stress matrices are standard
tools in rigidity theory; see, for example,
\cite{GortlerHealyThurston2010}.  We include the exact linear-algebra passage
needed here.

\begin{lemma}[all-nonzero Gale normalization]\label{lem:gale}
Assume $N-1\geq a$ and that the tensors in (3.1) are independent on a
nonempty Zariski-open set of $N\times a$ matrices.  There is a rational
$N\times a$ matrix $Z$, still having this independence property, such that
\[
 \rank Z=a,\qquad \sum_{i=1}^N z_i=0,
 \tag{4.1}
\]
and deleting any row of $Z$ leaves rank $a$.
\end{lemma}

\begin{proof}
Intersect the given independence locus with the nonempty open locus on which
every row deletion has rank $a$.  This intersection is a nonempty open set
defined over $\mathbb Q$, so choose a rational $Z$ in it.  Put
$W=\ker Z^{\mathsf T}$.  For each coordinate $i$, the row-deletion
condition implies that $W$ is not contained in the coordinate hyperplane
$c_i=0$: the other rows span $\K^a$, so there is a dependence involving
row $i$ with nonzero coefficient.  A finite union of proper hyperplanes
cannot cover $W$.  We may therefore choose a rational
$c=(c_1,\ldots,c_N)\in W$ with every $c_i\neq0$.  Replacing $z_i$ by
$c_i z_i$ gives (4.1), preserves every row-deletion rank, and multiplies
each edge tensor by the nonzero scalar $c_ic_j$.
\end{proof}

\begin{proposition}[stress restriction]\label{prop:stress}
Let $Z$ be as in Lemma~\ref{lem:gale}, suppose that the tensors indexed by
$E(D)$ are independent, and put $d=N-a-1$.  Then
$K_N-E(D)$ has a realization in $\mathbb Q^d$ whose rigidity matrix has
rank $dN-\binom{d+1}{2}$.
\end{proposition}

\begin{proof}
The kernel of $Z^{\mathsf T}$ has dimension $d+1$ and contains the all-ones
vector.  Choose an $N\times(d+1)$ matrix $[\mathbf1\ P]$ whose columns form
a basis of this kernel.  The rows of $P$ form an affinely spanning rational
realization $p$ in $\mathbb Q^d$.

For the complete framework $(K_N,p)$, write a stress as its symmetric stress
matrix $\Omega$: its off-diagonal entries are the negatives of the edge
stress coefficients and its diagonal entries make every row sum zero.  The
equilibrium equations are
\[
 \Omega\mathbf1=0,\qquad \Omega P=0.
\]
Because the columns of $Z$ span the orthogonal complement of the columns of
$[\mathbf1\ P]$, all complete-framework stress matrices are precisely
\[
 \Omega=Z\Psi Z^{\mathsf T},\qquad \Psi\in\Sym^2\K^a.
 \tag{4.2}
\]
In particular, restriction of (4.2) to a deleted edge $ij\in E(D)$ is the
linear functional
\[
 \Psi\longmapsto z_i^{\mathsf T}\Psi z_j.
\]
With respect to the trace pairing on symmetric matrices, the transpose of
this restriction map has columns $z_i\odot z_j$.  Their independence makes
the restriction map from the complete stress space onto
$\K^{E(D)}$ surjective.

Let $R_e$ denote the row of $R(K_N,p)$ indexed by $e$.  For any fixed
$e\in E(D)$, surjectivity supplies a complete stress whose coefficient on
$e$ is $1$ and whose coefficients on all other edges of $D$ are zero.  The
stress relation then writes $R_e$ as a linear combination of the rows
indexed by $E(K_N)\setminus E(D)$.  This holds for every deleted edge, so
\[
 \operatorname{rowspan}R(K_N-E(D),p)
 =\operatorname{rowspan}R(K_N,p).
\]
Since $p$ affinely spans $\K^d$, the latter rank is
$dN-\binom{d+1}{2}$.  The same is therefore true for the complement
framework.
\end{proof}

\begin{proof}[Proof of Theorem~\ref{thm:main}]
Let $a$ be minimal with $N\leq\binom{a+1}{2}$ and put $d=N-a-1$.
Minimality and $N\geq4$ imply $N-1\geq a$.  Proposition
\ref{prop:compression}, Lemma~\ref{lem:gale}, and Proposition
\ref{prop:stress} give a rational realization of $G$ with full rigidity
rank in dimension $d$.  A maximal minor that is nonzero at this realization
is a nonzero polynomial in the coordinates, so the generic rigidity matrix
has the same full rank.  Thus $G$ is generically $d$-rigid (with the
$d=0$ case covered by the standing convention).  Lemma~\ref{lem:maxwell}
gives $d=M(G)$, while Maxwell's necessary count gives the reverse bound
$\rig(G)\leq M(G)$.  Hence $\rig(G)=M(G)=N-a-1$.
\end{proof}

\section{Ancillary exact checks}

The proof above is uniform and does not rely on computation.  As adversarial
diagnostics, all feasible pairs $(s,t)$ with $4\leq3s+4t\leq28$ were also
checked at three deterministic integer realizations.  For each of the 41
pairs, exact elimination modulo $65521$ produced the full target rigidity
rank, and a separate verifier reconstructed the graph, placement, affine
span, selected maximal minor and its nonzero determinant.  A symbolic
transition check at $a=4$ also produced a determinant
$2q(27q+10)$ for ten edge tensors.  These calculations test small and sharp
regimes but are not evidence for untested orders.  For comparison, the same
methodological separation between an exact-arithmetic certificate and a
definition-level independent reconstruction is used in a different matrix
problem in \cite{ZengEtAl2026ExactPositivity}.

\section*{Computational and AI assistance}

OpenAI Codex assisted with proof discovery, exact finite and symbolic
experiments, literature screening, and preparation of this manuscript.  The
finite checks were performed in Python~3 with integer, finite-field, and
symbolic arithmetic; no floating-point rank decision enters the result.  A
separate fresh model context reconstructed the compression, rank-stratum,
Gale-normalization, and stress-surjectivity arguments and reran the supplied
exact checks.  This was an internal mathematical audit, not independent peer
review.  The universal conclusion rests on the deductive proof in
Sections~3--4.

\bibliographystyle{plain}
\bibliography{references}

\end{document}
